\documentclass[12pt,a4paper]{article}
\usepackage{graphicx}
\usepackage{amsmath}
\usepackage{amssymb}
\usepackage{amsthm}
\usepackage{bm}
\usepackage{hyperref}
\usepackage{natbib}
\usepackage{booktabs}
\usepackage{array}
\usepackage{geometry}
\geometry{margin=2.5cm}
\hypersetup{colorlinks=true, linkcolor=blue, citecolor=blue, urlcolor=blue}
\newcommand{\Jmax}{J_{\max}}
\newcommand{\Jreq}{J_{\mathrm{req}}}
\newcommand{\Jeff}{J_{\mathrm{eff}}}
\newcommand{\Om}{\Omega}
\newtheorem{theorem}{Theorem}
\newtheorem{proposition}{Proposition}
\newtheorem{definition}{Definition}
\newtheorem{lemma}{Lemma}
\newtheorem{corollary}{Corollary}
\newtheorem{criterion}{Criterion}

\title{\textbf{Paper BIO-07: Adaptive Flux Limitation in Oncology:\\ A CDFD Model of Constraint Escape, Resistance, and Boundary Failure}}
\author{Steve Bico Mujjabi\\ \small Independent Researcher}
\date{\today}

\begin{document}
\maketitle

\clearpage
\section*{Adaptive Flux Limitation in Oncology: A CDFD Model of Constraint Escape, Resistance, and Boundary Failure}

\begin{abstract}

This paper presents \textbf{Adaptive Flux Limitation (AFL)} as the Part III biology-and-medicine model inside Constraint-Driven Flux Dynamics (CDFD). CDFD supplies the flow-constraint-memory grammar: physiological flow $\Phi$, constraint $C$, surface responsiveness $S$, and structural memory $M_s$. The Runtime citation anchors the CDFL implementation, while clinical interpretation remains separate from software output and bedside validation.

Cancer is conventionally framed as a genetic disease of uncontrolled proliferation. While somatic mutations are essential, they do not explain the emergent systems-level behaviour of tumours: the Warburg effect, angiogenesis, immune evasion, metastasis, and multi-drug resistance. These are not independent hallmarks but coupled consequences of a single dynamical transition --- constraint escape.

We develop a deep AFL model of cancer in which tumour initiation is the crossing of a critical $\Psi_{s,\mathrm{growth}}$ threshold through constraint reduction, tumour progression is flux reprogramming to maximize throughput under the new constraint structure, and treatment failure is adaptive constraint re-establishment against each therapeutic intervention.

Key results: (1) carcinogenesis is a $C_{\mathrm{cell}} \to 0$ catastrophe triggered by oncogenic mutation of constraint-maintaining machinery (tumour suppressors, checkpoints, apoptosis); (2) the Warburg effect is a constraint \textit{restructuring} rather than a constraint \textit{removal} --- cancer trades oxidative efficiency for glycolytic speed under hypoxic constraint; (3) metastasis is triggered when local $C_{\mathrm{tissue}}$ exceeds the tumour's tolerance --- cells migrate to lower-constraint environments; (4) drug resistance is AFL adaptation: $\partial C_{\mathrm{drug}} / \partial t = \alpha |J_{\mathrm{drug}}| - \beta C_{\mathrm{drug}}$; (5) the tumour microenvironment (TME) is a constraint field that the tumour actively remodels.

\textbf{Status: Inferred}

\medskip
\noindent Within the extended framework of this series, biological networks are modeled as \emph{Adaptive Surfaces} that dynamically integrate physiological load and retain structural memory. The system's health is governed by the adaptive ratio $\Psi_s = (\Phi/C) \cdot S \cdot M_s$, where homeostasis ($\Psi_s \approx 1.0$) is maintained near the stable attractor ($S \cdot M_s \approx 1.0$), and disease emerges as surface dysregulation when maladaptive memory locks tissues into chronic constraint.

\end{abstract}


\paragraph{Greek-symbol convention.}
Unless a manuscript states a tighter equation-local meaning, Greek symbols are local CDFD variables or coefficients, not universal constants. Here $\Phi$ denotes driving flux, $\Psi_s$ denotes the adaptive operating ratio, $\Omega$ denotes a domain or state space, and $\Lambda$ denotes the Life Number where used. The coefficients $\alpha$, $\beta$, and $\gamma$ denote accumulation or reinforcement, relaxation or decay, and diffusion, coupling, or redistribution. The symbols $\delta$ and $\epsilon$ denote perturbation or tolerance scales, while $\eta$, $\lambda$, $\mu$, $\nu$, $\rho$, $\sigma$, $\tau$, and $\phi$ denote local efficiency, rate, baseline, frequency, density or correlation, noise or variance, time-scale, and phase or field terms. Subscripts restrict any coefficient to the named pathway, tissue, domain, or experiment.

\section*{Clinical Reader's Guide}
This oncology paper frames cancer as constraint escape and constraint re-engineering. A tumour is not only fast-growing tissue; it is a system that changes the rules under which growth, blood supply, immune visibility, redox handling, and therapy response are constrained. AFL names those changes as shifts in $\Phi$, $C$, $S$, and $M_s$.

The clinical reader should separate three levels. First, established cancer biology already describes metabolic reprogramming, angiogenesis, immune evasion, genomic instability, clonal selection, and resistance. Second, AFL groups these under a control architecture: the tumour lowers or bypasses constraints that would normally limit proliferation, then adapts when treatment imposes new constraints. Third, the runtime parasitic-drain diagnostic is only a toy analogy for a high-flow, low-system-responsiveness node.

The framework is useful only if it makes testable predictions about combination timing, adaptive resistance, metabolic vulnerability, or immune-boundary failure. It is not a substitute for oncology trials, tumour-specific biology, staging, molecular profiling, or standard-of-care treatment.


\vspace{0.5em}\noindent\fbox{\parbox{0.97\linewidth}{\small\textbf{AFL notation.} In this paper, $\Phi$ denotes the relevant physiological throughput, $C$ the operative biological constraint, $S$ surface responsiveness, and $M_s$ retained structural memory. When a dominant flow can be specified, $\Psi_s=(\Phi/C) \cdot S \cdot M_s$ denotes the operating ratio. Claim discipline: explicit assumptions, separated derivation vs. interpretation, and status labels.}}
\vspace{0.75em}

\section{Introduction}

Adaptive Flux Limitation (AFL) is the named model used in this paper; CDFD is the parent framework that supplies the variables, closure logic, and claim discipline. The biological or medical question is posed first as AFL, then interpreted as a CDFD model of flow, constraint, surface responsiveness, and structural memory.

Cancer is not simply rapid cell division. It is a living system that changes the constraints under which division, metabolism, immune visibility, tissue boundaries, and treatment response normally operate. Mutation is essential, but the clinical phenotype is architectural: proliferative signalling is amplified, checkpoints are bypassed, apoptosis is resisted, immune recognition is evaded, blood supply is recruited, and metastatic routes are opened.

AFL reads these hallmarks as changes in flow and constraint. Growth drive $\Phi_{\mathrm{growth}}$ rises, cell-cycle and apoptotic constraints fall, tissue and immune boundaries are remodelled, and structural memory $M_s$ accumulates through clonal selection and microenvironmental conditioning. The tumour is therefore not a single high-flow node; it is a constraint-engineering system embedded in a host.

This view is clinically useful only when it remains disciplined. It cannot replace histology, staging, molecular profiling, treatment guidelines, clinical trials, or patient-specific oncology judgement. Its value is to organize questions about adaptive resistance: which constraint did the treatment impose, which escape route did the tumour open, and which host surface paid the cost?

The runtime parasitic-drain analogy is deliberately limited. It is a toy diagnostic for a high-throughput node that can perturb surrounding responsiveness. It does not prove cancer biology. It suggests falsifiable ways to compare tumour load, immune boundary failure, metabolic state, and treatment timing.

\textbf{Status: Hypothesis}

\section{Carcinogenesis: The $C_{\mathrm{cell}} \to 0$ Catastrophe}

\subsection{Normal Cellular Constraint Architecture}

In a healthy epithelial cell, the effective constraint $C_{\mathrm{cell}}$ is maintained by a hierarchy of tumour suppressor mechanisms:
\begin{equation}
    C_{\mathrm{cell}} = C_{\mathrm{p53}} + C_{\mathrm{Rb}} + C_{\mathrm{PTEN}} + C_{\mathrm{APC}} + C_{\mathrm{contact}} + C_{\mathrm{apoptosis}} + \cdots
\end{equation}

Each component is a distinct constraint pathway. Loss of a single component reduces $C_{\mathrm{cell}}$ but may not cross the $\Psi_s > 1$ threshold alone. This is why carcinogenesis is multi-hit: typically 4--8 driver mutations are required.

\subsection{Multi-Hit Carcinogenesis as Progressive Constraint Erosion}

Each driver mutation removes one $C_i$ term:
\begin{align}
    \text{Hit 1 (e.g., APC loss):} & \quad C_{\mathrm{cell}} = C_{\mathrm{p53}} + C_{\mathrm{Rb}} + C_{\mathrm{PTEN}} + \cdots \quad \Psi_s < 1 \\
    \text{Hit 2 (e.g., KRAS):} & \quad \Phi_{\mathrm{growth}} \uparrow \quad \Psi_s \to 1 \\
    \text{Hit 3 (e.g., p53 loss):} & \quad C_{\mathrm{p53}} = 0 \quad \Psi_s > 1 \text{ (malignant)}
\end{align}

The AFL model predicts that the number of driver mutations required for malignant transformation is inversely proportional to the magnitude of each individual constraint. Tissues with higher baseline $C_{\mathrm{cell}}$ require more hits; tissues with lower baseline $C_{\mathrm{cell}}$ (e.g., rapidly dividing epithelium) are more cancer-prone.

\textbf{Status: Inferred}

\section{The Warburg Effect: Constraint Restructuring Under Hypoxia}

The Warburg effect --- preferential use of aerobic glycolysis even in the presence of oxygen --- is puzzling from an efficiency standpoint: glycolysis yields only 2 ATP per glucose vs 36--38 for OXPHOS.

The AFL explanation: the Warburg effect is not inefficiency but \textbf{constraint restructuring under spatial hypoxia}.

Solid tumours develop hypoxic cores as they outgrow their blood supply. Under hypoxia, the oxidative phosphorylation pathway faces a constraint:
\begin{equation}
    C_{\mathrm{OXPHOS}} \uparrow \text{ (oxygen limitation)} \Rightarrow J_{\mathrm{OXPHOS}} \downarrow
\end{equation}

Switching to glycolysis:
\begin{itemize}
    \item Removes the oxygen constraint ($C_{\mathrm{oxygen}}$ is irrelevant to glycolysis)
    \item Trades ATP yield for speed (glycolysis is $\sim$100$\times$ faster per enzyme molecule)
    \item Generates lactate that acidifies the TME, creating a new constraint on surrounding normal cells ($C_{\mathrm{normal}} \uparrow$ via acidosis)
\end{itemize}

In AFL terms: the Warburg effect is an adaptation that simultaneously reduces the tumour's own constraint ($C_{\mathrm{oxygen}}$ no longer limiting) and increases the constraint on the surrounding normal tissue (lactate acidification). It is constraint escape plus constraint imposition on competitors.

\textbf{Status: Inferred}

\section{Angiogenesis: Flux Amplification for Growing Constraint}

As a tumour grows, its intrinsic demand (nutrient flux $\Phi_{\mathrm{nutrients}}$) exceeds vascular supply. The gradient drives VEGF secretion, which stimulates angiogenesis.

In AFL terms:
\begin{equation}
    \Psi_{s,\mathrm{tumour}} = \frac{\Phi_{\mathrm{nutrients}}}{C_{\mathrm{vascular}}} < 1 \quad \text{(nutrient-limited)} \Rightarrow \text{VEGF} \uparrow \Rightarrow C_{\mathrm{vascular}} \downarrow
\end{equation}

Anti-VEGF therapy (bevacizumab) works by preventing this constraint reduction: it keeps $C_{\mathrm{vascular}}$ elevated, starving the tumour. But tumours adapt by upregulating alternative pro-angiogenic factors (FGF, PDGF), which is AFL constraint bypass: when one pathway for constraint reduction is blocked, another is activated.

\textbf{Anti-VEGF resistance AFL mechanism:}
\begin{equation}
    \frac{dC_{\mathrm{VEGF-resistance}}}{dt} = \alpha_{\mathrm{res}} |J_{\mathrm{anti-VEGF}}| - \beta_{\mathrm{res}} C_{\mathrm{VEGF-resistance}}
\end{equation}

Resistance emerges as the tumour adaptively increases its alternative angiogenic flux capacity in response to the drug-imposed constraint.

\textbf{Status: Inferred}

\section{The Tumour Microenvironment as a Constraint Field}

The TME --- comprising cancer-associated fibroblasts (CAFs), tumour-associated macrophages (TAMs), extracellular matrix (ECM), immunosuppressive T-regulatory cells, and myeloid-derived suppressor cells --- is not a passive bystander. The tumour actively remodels the TME to create a low-constraint environment.

\subsection{CAF-Mediated ECM Remodelling}

CAFs secrete matrix metalloproteinases (MMPs) that degrade collagen and fibronectin. In AFL terms: ECM degradation reduces the mechanical constraint $C_{\mathrm{ECM}}$, allowing tumour cells to expand with lower energy cost.
\begin{equation}
    C_{\mathrm{ECM}} \downarrow \Rightarrow \Phi_{\mathrm{growth}} / C_{\mathrm{ECM}} \uparrow \Rightarrow \text{expansion}
\end{equation}

\subsection{TAM Polarization: From M1 (Constraint) to M2 (Permissive)}

Tumour-associated macrophages can adopt M1 (anti-tumour, constraint-imposing) or M2 (pro-tumour, constraint-relaxing) phenotypes. The tumour secretes IL-4, IL-13, and TGF-$\beta$ to drive M2 polarization:
\begin{equation}
    \text{IL-4/IL-13} \Rightarrow \mathrm{TAM_{M2}} \Rightarrow C_{\mathrm{immune \to tumour}} \downarrow
\end{equation}

This is direct AFL constraint reduction by immunosuppressive remodelling of the TME.

\textbf{Status: Inferred}

\section{Immune Evasion: Checkpoint Exploitation as Constraint Reduction}

The immune system imposes constraint on tumour growth through cytotoxic T lymphocyte (CTL) activity:
\begin{equation}
    J_{\mathrm{immune}} = \frac{\Phi_{\mathrm{CTL}}}{ C_{\mathrm{immune}}} \quad \Rightarrow \quad \text{tumour cell killing}
\end{equation}

Tumour cells reduce $C_{\mathrm{immune \to tumour}}$ by upregulating checkpoint ligands:
\begin{itemize}
    \item \textbf{PD-L1}: binds PD-1 on CTLs, inducing T cell exhaustion ($\Phi_{\mathrm{CTL}} \downarrow$)
    \item \textbf{CTLA-4 pathway}: inhibits T cell activation at the priming stage
    \item \textbf{IDO pathway}: depletes tryptophan, creating metabolic constraint on T cells ($C_{\mathrm{T-cell}} \uparrow$)
\end{itemize}

Checkpoint inhibitors (pembrolizumab, nivolumab, ipilimumab) work by restoring $C_{\mathrm{immune \to tumour}}$ to normal levels. AFL prediction: checkpoint inhibitors should be most effective in tumours with highest PD-L1 expression (highest immune constraint reduction) and least effective in tumours with alternative immune evasion mechanisms not targeted by the drug.

\textbf{Status: Inferred}

\section{Metastasis: The Constraint Gradient Migration}

Metastasis occurs when tumour cells detach from the primary site and colonize distant organs. The AFL mechanism:

\begin{equation}
    C_{\mathrm{local}} > C_{\mathrm{tumour\ tolerance}} \Rightarrow \text{migration signal} \Rightarrow J_{\mathrm{cells}} \to \text{lower-constraint sites}
\end{equation}

The local constraint $C_{\mathrm{local}}$ rises in the primary tumour as:
\begin{itemize}
    \item Hypoxia increases (oxygen constraint)
    \item Lactate acidification increases (pH constraint)
    \item Immune pressure increases (immune constraint)
    \item Nutrient competition increases (metabolic constraint)
\end{itemize}

When $C_{\mathrm{local}}$ exceeds the tumour's adaptation capacity, cells undergo epithelial-mesenchymal transition (EMT) --- a phenotypic switch that reduces cellular adhesion (lowers $C_{\mathrm{contact}}$) and enables motility.

\textbf{Metastatic organotropism} (why specific cancers colonize specific organs) is predicted by the AFL model as a matching of tumour constraint profile to target organ constraint profile: breast cancer metastasizes to bone because bone ECM provides a mechanical constraint environment similar to the breast stroma; colorectal cancer metastasizes to liver because portal venous flow delivers cells to a high-nutrient, low-immune constraint environment.

\textbf{Status: Inferred}

\section{Drug Resistance: AFL Adaptation to Therapeutic Constraints}

\subsection{The Universal Resistance Mechanism}

Every effective cancer therapy imposes a new constraint $C_{\mathrm{drug}}$ on the tumour. The tumour's response follows the AFL adaptation equation:
\begin{equation}
    \frac{dC_{\mathrm{drug-resistance}}}{dt} = \alpha_{\mathrm{res}} |J_{\mathrm{drug}}| - \beta_{\mathrm{res}} C_{\mathrm{drug-resistance}}
\end{equation}

As drug concentration $J_{\mathrm{drug}}$ rises, resistance $C_{\mathrm{drug-resistance}}$ rises. At steady state:
\begin{equation}
    C_{\mathrm{drug-resistance}}^* = \frac{\alpha_{\mathrm{res}}}{\beta_{\mathrm{res}}} |J_{\mathrm{drug}}|
\end{equation}

Higher drug concentration drives faster resistance evolution. This is the AFL explanation of the ``dose-escalation paradox'': increasing dose accelerates resistance rather than cure.

\subsection{Resistance Mechanisms as AFL Bypass Strategies}

Specific resistance mechanisms map to AFL strategies:
\begin{itemize}
    \item \textbf{Efflux pumps (P-gp)}: actively reduce intracellular drug concentration ($\Phi_{\mathrm{drug,intracellular}} \downarrow$)
    \item \textbf{Target mutation}: reduces drug-target binding, raising the effective $C_{\mathrm{drug}}$ threshold for effect
    \item \textbf{Pathway redundancy}: activates alternative signalling routes ($C_{\mathrm{bypass}} \downarrow$ via alternative pathway)
    \item \textbf{Senescence bypass}: cancer cells that would normally undergo therapy-induced senescence instead bypass the checkpoint ($C_{\mathrm{senescence}} \downarrow$)
\end{itemize}

\textbf{Status: Inferred}

\subsection{Combination Therapy as Multi-Constraint Imposition}

AFL predicts: effective treatment must impose constraints on multiple pathways simultaneously, because single-pathway constraint increase drives bypass evolution. The mathematical argument:

If the tumour has $n$ bypass pathways with resistance rates $\alpha_1, \ldots, \alpha_n$, effective resistance develops at rate $\min_i \alpha_i$ under monotherapy but requires \textit{all} bypass pathways to evolve simultaneously under combination therapy --- a much slower process.

This predicts the superiority of combination regimens and the clinical observation that resistance to combination chemotherapy develops far more slowly than resistance to single agents.

\textbf{Status: Inferred}

\section{Flux--Constraint Regime Table}

\begin{center}
\begin{tabular}{lll}
\toprule
$\Psi_{s,\mathrm{tumour}}$ range & State & Clinical correlate \\
\midrule
$< 1$ (sustained) & Controlled & Dormant micrometastasis, treatment-suppressed \\
$\approx 1$ & Pre-malignant & CIN, dysplasia, carcinoma in situ \\
$> 1$ (acute) & Active growth & Invasive carcinoma \\
$\gg 1$ & Aggressive & High-grade, rapidly dividing tumour \\
After therapy: $\to 0$ & Treatment response & Complete or partial response \\
Rising after nadir & Resistance & Progressive disease, recurrence \\
\bottomrule
\end{tabular}
\end{center}

\textbf{Status: Inferred}

\section{Predictions}

\begin{enumerate}
    \item \textbf{Multi-hit threshold varies by tissue}: Cancer incidence rate across tissue types should correlate with the average number of tumour suppressor genes expressed per cell (higher expression = higher baseline $C_{\mathrm{cell}}$ = more hits required = lower incidence).
    \item \textbf{Warburg switch timing}: Aerobic glycolysis should activate when intratumoral $\mathrm{pO_2}$ falls below a threshold, not at a fixed tumour size --- because the trigger is oxygen constraint, not size per se.
    \item \textbf{Metastasis timing prediction}: The time from primary diagnosis to first metastasis should correlate with the rate of TME constraint increase in the primary site, measurable via CAF density and ECM stiffness on biopsy.
    \item \textbf{Combination therapy resistance rate}: The time to resistance under combination therapy should scale with $\sum_i \alpha_i^{-1}$ (sum of inverse resistance rates), which is measurable from single-agent resistance timecourses.
    \item \textbf{Checkpoint inhibitor response prediction}: Response to anti-PD-1 should correlate with the ratio of PD-L1 expression to alternative immune evasion pathway activity, not with PD-L1 alone.
\end{enumerate}

\textbf{Status: Open}

\section{Falsifiability}

The AFL oncology model would be falsified if:
\begin{itemize}
    \item Cancer incidence is uncorrelated with tumour suppressor gene count across tissues.
    \item The Warburg effect is present in well-oxygenated tumour cells without alternative metabolic constraint.
    \item Metastasis occurs at fixed tumour sizes regardless of TME constraint dynamics.
    \item Combination therapy resistance develops as rapidly as monotherapy resistance.
    \item Checkpoint inhibitor response is independent of alternative immune evasion pathway activity.
\end{itemize}

\textbf{Status: Open}

\section{Limitations}

\begin{itemize}
    \item Cancer heterogeneity (multiple subclonal populations with different $C_{\mathrm{cell}}$) is not fully captured by a single-constraint scalar model.
    \item The model does not distinguish between driver and passenger mutations in terms of constraint contribution.
    \item Immunotherapy combinations (anti-PD1 + anti-CTLA4) involve constraint interactions between two immune pathways not fully captured here.
    \item Liquid tumours (leukaemia, lymphoma) have fundamentally different spatial constraint architecture than solid tumours --- this model primarily addresses solid tumours.
\end{itemize}

\textbf{Status: Open}


\section{Deep Analysis: Resistance ODE: Global Stability and Compensatory Mutations}
\noindent\textit{Detailed derivation placeholder retained for future expansion in this preprint release.}


\section{Biological Mujjabi Laws and Diagnostics}

\subsection{The Mujjabi Parasitic-Drain Law \cite{MujjabiPartII2026} (Biology)}
Inspired by the Part II error-threshold discussion, the \textbf{Mujjabi Parasitic-Drain Law \cite{MujjabiPartII2026}} proposes that localized high-flow nodes with low systemic responsiveness ($S \ll 1$) can impose disproportionate constraint on the surrounding adaptive surface. Candidate biological analogies include tumours, viral replication centers, biofilms, and chronic inflammatory foci.

\subsubsection{Runtime Diagnostic: Parasitic Drain}
Release-local runtime diagnostics in \texttt{outputs/paper\_07/} introduce a high-flow, low-responsiveness node and track how it perturbs the surrounding operating ratio. This is a toy analogy for tumour-like or chronic-focus behavior. It does not prove a cancer mechanism, and it should not be read as a substitute for tumour biology, staging, molecular profiling, or clinical oncology evidence.

\section{Clinical Mapping of Symbols}

\begin{center}
\begin{tabular}{p{1.5cm}p{3.0cm}p{3.5cm}p{4.5cm}}
\toprule
Symbol & Oncology meaning & Measurement proxy & Invalidation condition \\
\midrule
$\Phi$ & Proliferative flux / growth drive & Ki-67, ctDNA fraction, tumour growth rate & If growth rate is uncorrelated with $\Phi$ proxy, mapping fails \\
$C$ & Constraint on proliferation (apoptotic, immune) & p53 status, immune infiltration score, drug concentration & If constraint escape is explained by a single mutation without adaptive dynamics, AFL adds nothing \\
$M_s$ & Epigenetic / phenotypic lock & DNA methylation pattern, chromatin accessibility, persister cell fraction & If $M_s$ reverses completely after drug removal, locked-state framing is not needed \\
$\Psi_s$ & Tumour-host operating state & Computed from tumour/host constraint ratio; no direct assay & If $\Psi_s$ trajectory matches a simple PD-L1 score, framework adds no information \\
\bottomrule
\end{tabular}
\end{center}

\textbf{Status: Hypothesis}

\section{Known Medicine Baseline}

\begin{itemize}
    \item Oncology evidence base: NCI TCGA/CPTAC genomics/proteomics, established tumour staging, approved targeted therapies and immunotherapy — all based on clinical trial evidence.
    \item \textbf{FDA Project Optimus}: dose optimization as efficacy plus safety/tolerability, not maximum tolerated dose alone. AFL ``axis orthogonality'' in combination therapy is consistent with this principle but does not replace trial-based dose-ranging.
    \item Tumour microenvironment (hypoxia, angiogenesis, immune exclusion, stromal stiffness): established oncology domains with active research; AFL constraint language is a research reframing.
    \item When AFL framing conflicts with oncology guidelines or trial-based prescribing, the trial evidence takes precedence.
\end{itemize}

\section{Clinical Safety Boundary}

This paper is a research framework, not a treatment guide.
\begin{itemize}
    \item Cancer treatment decisions (surgery, chemotherapy, immunotherapy, radiotherapy, targeted therapy) require oncology specialist assessment, molecular profiling, staging, and multidisciplinary tumour board review.
    \item AFL does not provide predictive biomarkers, response prediction, or dose recommendations.
    \item Tumour escape predictions are research hypotheses requiring organoid, ctDNA, or spatial transcriptomics validation.
\end{itemize}

\section{Runtime Diagnostic}

Release-local script: \texttt{supplementary/AFL\_Biological\_Mujjabi\_Tests.py} (parasitic drain section). Outputs in \texttt{outputs/paper\_07/}: \texttt{parasitic\_drain.png}, \texttt{parasitic\_drain\_timeseries.csv}, \texttt{summary.json}. All outputs are toy diagnostics; no clinical inference should be drawn.

\section{Conclusion}

Cancer is a constraint escape system. AFL unifies carcinogenesis (multi-hit $C$ erosion), tumour progression (Warburg/angiogenesis flux reprogramming), drug resistance (AFL adaptation), and immune evasion (checkpoint-mediated constraint reduction) under a single framework. Combination therapy and multi-constraint biomarkers follow as research predictions requiring NCI TCGA/CPTAC validation and FDA Project Optimus-aligned trial design.

\textbf{Status: Inferred}

\section{Reproducibility Note}

Release-local script: \texttt{supplementary/AFL\_Biological\_Mujjabi\_Tests.py} (NumPy + matplotlib; runtime $< 2$\,s).

\textbf{Status: Derived}
\clearpage
\section*{Adaptive Flux Limitation in Oncology: Cancer as a Constraint-Escape and Flux-Reprogramming System}


\bibliographystyle{plainnat}
\bibliography{afl_biology_references}

\end{document}
